% Ad Familiares 7.20 --- headnote
% ad-familiares-07-20, 20 July 44 BC, Velia

\textit{Cicero to C. Trebatius Testa, written from Velia on 20 July
44 BC, as Cicero made his way down the Italian coast on the abortive
voyage toward Greece that the south wind would soon turn back at
Leucopetra. Velia was Trebatius's birthplace, and Cicero spends the
opening on the town's affection for its absent jurist --- and on
Rufio, a member of Trebatius's household whom Trebatius has hauled
away to work on building projects in Rome, leaving the Velians
short of him.}

\textit{The substantive note --- delivered as a friend's advice in
unsettled times --- is that Trebatius should not sell or abandon the
paternal estate, with its noble river Halaes and its Papirian house
(though the lotus there might do with cutting down, for both view
and prudence). A house in a healthy, remote, friendly place is the
right thing to have, ``especially in these times.'' The closing
section turns the jurist's customary teasing back on Cicero himself:
he has carried off a book of Nico the physician's \textit{On
Overeating} from Nico's pupil Sex. Fadius, and is, of course,
ideally apt as a pupil in that discipline --- though their friend
Bassus had been hiding it from him, while apparently sharing it
with Trebatius.}
